\documentclass[aspectratio=169,11pt]{beamer}
\usepackage[T1]{fontenc}
\usepackage{lmodern}
\usepackage{amsmath,amssymb}
\definecolor{ink}{HTML}{173247}
\definecolor{accent}{HTML}{007E87}
\definecolor{paper}{HTML}{FAFAF7}
\setbeamercolor{normal text}{fg=ink,bg=paper}
\setbeamercolor{frametitle}{fg=ink}
\setbeamercolor{structure}{fg=accent}
\setbeamerfont{frametitle}{size=\Large,series=\bfseries}
\setbeamertemplate{navigation symbols}{}
\setbeamertemplate{footline}{\hfill\usebeamercolor[fg]{structure}\small\insertframenumber\hspace{7mm}\vspace{4mm}}
\setbeamersize{text margin left=10mm,text margin right=10mm}
\setbeamertemplate{itemize items}[circle]
\newcommand{\key}[1]{{\color{accent}#1}}
\begin{document}

% Source: lambda_calculus_intuitive_explanation.pdf, introduction and section 7.
\begin{frame}
\vspace{7mm}
{\Huge\bfseries An Intuitive Introduction\\[3mm]to Lambda Calculus\par}
\vspace{9mm}
{\Large Computation with functions\par}
\vspace{12mm}
{\normalsize CS413\par}
\vspace{3mm}
{\small Based on \emph{An Intuitive Introduction to Lambda Calculus}}
\end{frame}

% Source: section 1, p. 1.
\begin{frame}{Anonymous functions}
An ordinary function has a name:
\[f(x)=x+1\]
Lambda notation writes the function itself:
\[\key{\lambda x.\,x+1}\]
Read it as ``take $x$ and return $x+1$.''
\vspace{4mm}

Python uses the same idea: \quad\texttt{lambda x: x + 1}
\end{frame}

% Source: section 2, pp. 1--2. Arithmetic distinction made explicit.
\begin{frame}{Application and beta reduction}
Applying a function supplies an argument:
\[\key{(\lambda x.\,x+1)\;5}\]
\begin{align*}
(\lambda x.\,x+1)\;5 &\longrightarrow_{\beta} 5+1 \\
&=6
\end{align*}
Beta reduction substitutes the argument for the parameter.
\vspace{5mm}

{\small Here, numbers and addition are familiar teaching notation. The last step is arithmetic.}
\end{frame}

% Source: section 3, p. 2.
\begin{frame}{The entire syntax of pure lambda calculus}
\begin{align*}
x &\qquad \text{variable}\\[3mm]
\lambda x.\,M &\qquad \text{abstraction (a function)}\\[3mm]
M\;N &\qquad \text{application}
\end{align*}
\vspace{3mm}
$M$ and $N$ stand for expressions.
\vspace{5mm}

Pure lambda calculus has no built-in numbers, addition, Booleans, or loops.
\end{frame}

% Source: section 4, pp. 2--3. Parenthesization is an explanatory addition.
\begin{frame}{Functions as arguments and results}
\[\key{\lambda f.\,\lambda x.\,f\;x}\]
First take $f$. Return a function that takes $x$ and applies $f$ to it.
\vspace{3mm}
\begin{align*}
((\lambda f.\,\lambda x.\,f\;x)\;(\lambda y.\,y+1))\;5
&\longrightarrow_{\beta}(\lambda x.\,(\lambda y.\,y+1)\;x)\;5\\
&\longrightarrow_{\beta}(\lambda y.\,y+1)\;5\\
&\longrightarrow_{\beta}5+1=6
\end{align*}
{\small Application groups to the left: $M\;N\;P$ means $(M\;N)\;P$.}
\end{frame}

% Source: section 5, p. 3.
\begin{frame}{Church numerals: numbers as repetition}
A numeral tells us how many times to apply a function.
\begin{align*}
\mathbf{0}&\equiv\lambda f.\,\lambda x.\,x\\[2mm]
\mathbf{1}&\equiv\lambda f.\,\lambda x.\,f\;x\\[2mm]
\mathbf{2}&\equiv\lambda f.\,\lambda x.\,f\;(f\;x)\\[2mm]
\mathbf{3}&\equiv\lambda f.\,\lambda x.\,f\;(f\;(f\;x))
\end{align*}
\key{The function's behavior encodes the number.}
\end{frame}

% Source: section 5, p. 3. Exercise is derived from the same definition.
\begin{frame}{Using a Church numeral}
Let $\mathrm{step}=\lambda n.\,n+10$.
\begin{align*}
\mathbf{3}\;\mathrm{step}\;0
&\longrightarrow_{\beta}^{*}\mathrm{step}\;(\mathrm{step}\;(\mathrm{step}\;0))\\[2mm]
&=\mathrm{step}\;(\mathrm{step}\;10)\\[2mm]
&=\mathrm{step}\;20\\[2mm]
&=30
\end{align*}
{\small $\longrightarrow_{\beta}^{*}$ denotes zero or more beta-reduction steps. Arithmetic remains teaching notation.}
\vspace{4mm}

\key{Quick check:} What does $\mathbf{2}\;\mathrm{step}\;5$ produce?
\end{frame}

% Source: section 6, p. 4.
\begin{frame}{Booleans as choices}
\begin{align*}
\mathrm{TRUE}&\equiv\lambda x.\,\lambda y.\,x\\[3mm]
\mathrm{FALSE}&\equiv\lambda x.\,\lambda y.\,y
\end{align*}
TRUE chooses the first argument.\\[2mm]
FALSE chooses the second argument.
\vspace{3mm}
\begin{align*}
\mathrm{TRUE}\;A\;B&\longrightarrow_{\beta}^{*}A\\
\mathrm{FALSE}\;A\;B&\longrightarrow_{\beta}^{*}B
\end{align*}
\end{frame}

% Source: section 6, p. 4.
\begin{frame}{A conditional built from functions}
\[\key{\mathrm{IF}\equiv\lambda b.\,\lambda x.\,\lambda y.\,b\;x\;y}\]
The Boolean $b$ selects one of the two branches.
\begin{align*}
\mathrm{IF}\;\mathrm{TRUE}\;A\;B
&\longrightarrow_{\beta}^{*}\mathrm{TRUE}\;A\;B
\longrightarrow_{\beta}^{*}A\\[4mm]
\mathrm{IF}\;\mathrm{FALSE}\;A\;B
&\longrightarrow_{\beta}^{*}\mathrm{FALSE}\;A\;B
\longrightarrow_{\beta}^{*}B
\end{align*}
\vspace{4mm}
Boolean choice needs only abstraction and application.
\end{frame}

% Source: sections 7--8, pp. 4--5. Free/bound example and exercise answer added.
\begin{frame}{Computational power and the next step}
Variables, abstraction, and application suffice to express arbitrary computation.
\vspace{4mm}

Functions can encode numbers, Booleans, data structures, and recursion.
\vspace{5mm}

\key{Next: precise substitution and variable binding}
\[\lambda x.\,x\;y\qquad\text{$x$ is bound, $y$ is free}\]
{\small Substitution must preserve the meaning of free variables.}
\vspace{4mm}

{\small Quick-check answer: $\mathbf{2}\;\mathrm{step}\;5=25$.}
\end{frame}
\end{document}
